From \parencite{sarker2021machine_learning_algorithms_real_world_applications_research_directions} :

"Supervised learning is typically the task of machine learning to learn a function that maps an input to an output based on sample input-output pairs. It uses labeled training data and a collection of training examples to infer a function [...] The most common supervised tasks are 'classification' that separates the data, and 'regression' that fits the data."

"Unsupervised learning analyzes unlabeled datasets without the need for human interference, i.e., a data-driven process. This is widely used for extracting generative features, identifying meaningful trends and structures, groupings in results, and exploratory purposes. The most common unsupervised learning tasks are clustering, density estimation, feature learning, dimensionality reduction, finding association rules, anomaly detection, etc."

"Semi-supervised learning can be defined as a hybridization of the above-mentioned supervised and unsupervised methods, as it operates on both labeled and unlabeled data [...] In the real world, labeled data could be rare in several contexts, and unlabeled data are numerous, where semi-supervised learning is useful."

"Reinforcement learning is a type of machine learning algorithm that enables software agents and machines to automatically evaluate the optimal behavior in a particular context or environment to improve its efficiency, i.e., an environment-driven approach. This type of learning is based on reward or penalty, and its ultimate goal is to use insights obtained from environmental activists to take action to increase the reward or minimize the risk."

From \parencite[p.~3]{bishop2006pattern_recognition_machine_learning} :

"Applications in which the training data comprises examples of the input vectors along with their corresponding target vectors are known as supervised learning problems. Cases such as the digit recognition example, in which the aim is to assign each input vector to one of a finite number of discrete categories, are called classification problems. If the desired output consists of one or more continuous variables, then the task is called regression."

"In other pattern recognition problems, the training data consists of a set of input vectors x without any corresponding target values. The goal in such unsupervised learning problems may be to discover groups of similar examples within the data, where it is called clustering, or to determine the distribution of data within the input space, known as density estimation, or to project the data from a high-dimensional space down to two or three dimensions for the purpose of visualization."

"The technique of reinforcement learning is concerned with the problem of finding suitable actions to take in a given situation in order to maximize a reward. Here the learning algorithm is not given examples of optimal outputs, in contrast to supervised learning, but must instead discover them by a process of trial and error."

From \parencite{Goodfellow-et-al-2016} :

"Applications in which the training data comprises examples of the input vectors along with their corresponding target vectors are known as supervised learning problems. Cases such as the digit recognition example, in which the aim is to assign each input vector to one of a finite number of discrete categories, are called classification problems. If the desired output consists of one or more continuous variables, then the task is called regression."

"In other pattern recognition problems, the training data consists of a set of input vectors x without any corresponding target values. The goal in such unsupervised learning problems may be to discover groups of similar examples within the data, where it is called clustering, or to determine the distribution of data within the input space, known as density estimation, or to project the data from a high-dimensional space down to two or three dimensions for the purpose of visualization."

"The technique of reinforcement learning is concerned with the problem of finding suitable actions to take in a given situation in order to maximize a reward. Here the learning algorithm is not given examples of optimal outputs, in contrast to supervised learning, but must instead discover them by a process of trial and error."

"Roughly speaking, unsupervised learning involves observing several examples of a random vector x, and attempting to implicitly or explicitly learn the probability distribution p(x) [...] while supervised learning involves observing several examples of a random vector x and an associated value or vector y, and learning to predict y from x."

"Unsupervised learning and supervised learning are not formally defined terms. The lines between them are often blurred."

From \parencite{mitchell1997machine_learning} :

"A computer program is said to learn from experience E with respect to some class of tasks T and performance measure P, if its performance at tasks in T, as measured by P, improves with experience E."

